\documentclass[12pt,a4paper]{article}

\usepackage[spanish]{babel}
\usepackage[utf8]{inputenc}
\usepackage[T1]{fontenc}
\usepackage{geometry}
\usepackage{booktabs}
\usepackage{enumitem}
\usepackage{xcolor}
\usepackage{hyperref}
\usepackage{listings}

\geometry{margin=2.5cm}

\title{Documentación de despliegue}
\author{EventConnect}
\date{\today}

\begin{document}

\maketitle

\section{Introducción}

\noindent
Este documento describe el proceso de despliegue de la aplicación \textbf{EventConnect}.
La aplicación está formada por tres partes principales: el frontend, el backend y la
base de datos. El frontend y el backend se encuentran desplegados en \textbf{Render},
mientras que la base de datos se gestiona mediante \textbf{MongoDB Atlas}. \\

\noindent
El objetivo de esta documentación es explicar cómo está desplegada la aplicación, qué
servicios intervienen, cómo se realiza el arranque inicial y cómo se cargan los datos
necesarios para que el sistema funcione correctamente.

\section{Servicios utilizados}

La aplicación utiliza los siguientes servicios para su despliegue:

\begin{center}
\begin{tabular}{ll}
  \toprule
  \textbf{Componente} & \textbf{Servicio utilizado} \\
  \midrule
  Frontend & Render, como sitio estático \\
  Backend & Render, como servicio web \\
  Base de datos & MongoDB Atlas \\
  Repositorios & GitHub \\
  \bottomrule
\end{tabular}
\end{center}

\noindent
El despliegue se ha organizado separando el frontend y el backend en servicios
independientes.
\section{Despliegue de la aplicación}

\noindent
La aplicación está desplegada utilizando \textbf{Render} para el frontend y el backend,
y \textbf{MongoDB Atlas} como base de datos en la nube. La base de datos se encuentra
en un cluster de MongoDB Atlas, desde donde se obtiene la URI de conexión. Para poder
acceder a ella se ha creado un usuario con contraseña y se ha permitido el acceso desde
cualquier IP mediante la regla \texttt{0.0.0.0/0}, ya que Render no proporciona una IP
fija en este tipo de despliegues. La URI de conexión se configura en Render como
variable de entorno, evitando incluir credenciales directamente en el código. \\

\noindent
El backend está desplegado en Render como un \textbf{servicio web}. Este servicio está
conectado al repositorio del backend en GitHub y utiliza una rama específica de
despliegue. Cuando se suben cambios a esa rama, Render permite desplegar la nueva
versión del backend. En este servicio se configuran las variables de entorno necesarias,
como la conexión a MongoDB, las claves de autenticación y las claves de servicios
externos. \\

\noindent
El frontend también está desplegado en Render, pero como \textbf{sitio estático}. Al
igual que el backend, está conectado a su repositorio de GitHub y utiliza una rama de
despliegue. Render construye la aplicación Angular y publica los archivos generados
para que puedan ser accedidos desde el navegador. En este caso no se han configurado
variables de entorno adicionales, ya que el frontend no necesita ninguna. \\

\section{Carga inicial de datos} 

\noindent
Para que la aplicación tenga datos disponibles, se utilizan procesos de carga inicial y
sincronización. La aplicación cuenta con \textbf{seeders}, que permiten cargar datos
iniciales en la base de datos cuando sea necesario. \\

\noindent
Además, el backend incluye funcionalidades para importar eventos desde la API de Datos
Abiertos del Ayuntamiento de Zaragoza. Esta importación permite guardar eventos reales
en MongoDB y utilizarlos posteriormente desde el frontend y desde las funcionalidades
de IA.

\section{URLs de acceso}

Las URLs públicas de acceso son las siguientes:

\begin{center}
\begin{tabular}{ll}
  \toprule
  \textbf{Recurso} & \textbf{URL} \\
  \midrule
  Frontend & \textit{https://frontend-c9bg.onrender.com/} \\
  API / Swagger & \textit{https://backend-hi6i.onrender.com/api/docs/} \\
  \bottomrule
\end{tabular}
\end{center}

\end{document}
